\subsection{Runtime lower bound}

Our goal here is to present a reduction showing that the algorithm
we have given for \SF parameterized by vertex cover is essentially optimal,
assuming the ETH. Recall that the ETH is the hypothesis that \textsc{3-SAT} on
instances with $n$ variables cannot be solved in time $2^{o(n)}$. We will give
a reduction that given a \textsc{3-SAT} instance $\phi$, produces an equivalent
\SF instance with vertex cover at most $O(n/\log n)$. We stress that our
reduction works even for unweighted instances.


\begin{theorem} If there exists an algorithm which, given an unweighted \SF
instance on~$n$ vertices with vertex cover $k$, finds an optimal solution in
time $2^{o(k\log k)}n^{O(1)}$, then the ETH is false. \end{theorem}

\begin{proof}

We present a reduction from \textsc{3-SAT}. Before we proceed, we would like to
add to our formula the requirement that the variable set comes partitioned into
three sets in a way that each clause contains at most one variable from each
set. It is not hard to show that this does not affect the complexity of the
instance much, as we demonstrate in the following claim.

\begin{claim}\label{claim:tsat3} %
Suppose that there exists an algorithm that takes as input a
\textsc{3-SAT} instance $\phi$ on $n$ variables and a partition of the
variables into three sets of equal size, such that each clause contains at most
one variable from each set, and decides if $\phi$ is satisfiable in time
$2^{o(n)}$. Then, the ETH is false. \end{claim}

%
\begin{proof}%
Suppose we start with an arbitrary \textsc{3-SAT} formula $\psi$
on $n$ variables $x_1,\ldots,x_n$. Under the ETH, it should be impossible to
decide if $\psi$ is satisfiable in time $2^{o(n)}$. We will edit $\psi$ to
produce the partition of the variables into three sets. For each variable
$x_i$, we introduce two new variables $x_i', x_i''$, and add to the formula the
clauses $(x_i\to x_i')\land (x_i'\to x_i'')\land (x_i''\to x_i)$.  The
variables of $\psi$ can now be partitioned into three sets $X =
\{x_1,\ldots,x_n\}$, $X'=\{ x_1',\ldots,x_n'\}$, and $X''=\{
x_1'',\ldots,x_n''\}$. Furthermore, because of the clauses we added it is not
hard to see that in any satisfying assignment $x_i, x_i'$, and $x_i''$ must be
given the same value. We then repeat the following: as long as there exists a
clause that contains more than one variable from $X$, arbitrarily pick a
literal of this clause that contains $x_i\in X$ and replace in it $x_i$ by
$x_i'$ or $x_i''$, in a way that the clause contains at most one variable from
each group. The new formula we have constructed in this way is equisatisfiable
to $\psi$, has $n'=O(n)$ variables and $O(n+m)$ clauses, and its variables are
partitioned into three sets so that each clause contains at most one variable
from each set. Therefore, the new formula cannot be solved in time $2^{o(n')}$
under the ETH.  \end{proof}
    
%

In the remainder we will then assume that we are given a formula $\phi$ on $3n$
variables which are partitioned into three sets of size $n$ as specified by the
previous claim. Without loss of generality, suppose that $n$ is a power of $4$
(this can be achieved by adding dummy variables). Note that this ensures that
$\frac{\log n}{2}$ and $\sqrt{n}$ are both integers.

We construct an equivalent instance of \SF as follows. Let
$L=\lceil\frac{n}{\log^2n}\rceil$. We  begin by constructing $i$ choice
gadgets, i.e., for $i\in \{1,\ldots,3\log n\}$ we make:

\begin{itemize}

\item $2L$ left vertices, labeled $\ell^i_j$, for $j\in\{0,\ldots,2L-1\}$.

\item $2L$ right vertices, labeled $r^i_j$, for $j\in\{0,\ldots,2L-1\}$.

\item $\sqrt{n}$ middle vertices, labeled $m^i_j$, for
$j\in\{0,\ldots,\sqrt{n}-1\}$. 

\item We connect all middle vertices to all left and right vertices, that is,
for all $j\in\{0,\ldots,2L-1\}$ and $j'\in\{0,\sqrt{n}-1\}$ we connect
$\ell^i_j$ and $r^i_j$ to $m^i_{j'}$.

\item For each $j\in\{0,\ldots,2L-1\}$ we add a demand from $\ell^i_j$ to
$r^i_j$.

\end{itemize}

Notice that the graph we have constructed so far contains $3\log n$ choice
gadgets, each of which has $4L+\sqrt{n}=O(n/\log^2n)$ vertices, so the graph at
the moment contains $O(n/\log n)$ vertices in total.

Before we proceed, let $X=X_a\cup X_b\cup X_c$ be the set of $3n$ variables of
$\phi$ that was given to us partitioned into three sets of size $n$. We
partition $X$ into $3\log n$ groups $X_1,\ldots,X_{3\log n}$ in a way that (i)
$|X_i|\le \lceil n/\log n\rceil$ for all $i\in\{1,\ldots,\log n\}$ and (ii) for
all $i\in\{1,\ldots,\log n\}$ we have $X_i$ is contained in one of
$X_a,X_b,X_c$. This can be done by taking the $n$ variables of $X_a$ and
partitioning them arbitrarily into groups $X_1,\ldots,X_{\log n}$ of size as
equal as possible (therefore at most $\lceil n/\log n\rceil$), and we proceed
similarly for $X_b,X_c$. Rename the variables of $\phi$ so that for each $i$ we
have that $X_i=\{x_{(i,0)},\ldots,x_{(i,\lceil n/\log n \rceil-1)}\}$.

To give some intuition, we will now say that, for $i\in\{1,\ldots, 3\log n\}$,
the choice gadget $i$ represents the variables of the set $X_i$. In particular,
for each $j\in\{0,\ldots 2L-1\}$, we will say that the way that the demand
$\ell^i_j\to r^i_j$ was satisfied encodes the assignment to the $\frac{\log
n}{2}$ variables $\{x_{(i,\frac{j\log n}{2})},\ldots,x_{(i,\frac{(j+1)\log
n}{2}-1)}\}$. More precisely, in our intended solution the demand $\ell^i_j\to
r^i_j$ is satisfied by connecting both terminals to a common middle vertex
$m^i_{j'}$. We can infer the assignment to the $\frac{\log n}{2}$ variables
this represents simply by writing down the binary representation of $j'$, which
is a number between $0$ and $\sqrt{n}-1$, hence a number with $\frac{\log
n}{2}$ bits. Note that this way we represent $2L \cdot \frac{\log n}{2}\ge
\lceil\frac{n}{\log n}\rceil$ variables, that is, we can represent the
assignment to all the variables of the group.

Armed with this intuition, we can now complete our construction. For each
clause $c$ we construct two new vertices, $c_1,c_2$ and add a demand from $c_1$
to $c_2$.  For each literal contained in $c$, suppose that the literal involves
the variable $x_{(i,\frac{j\log n}{2}+\alpha)}$ for $i\in\{1,\ldots,3\log n\}$,
$j\in \{0,\ldots,2L-1\}$, $\alpha\in\{0,\ldots,\frac{\log n}{2}-1\}$. We then
connect $c_1$ to $\ell^i_j$. Furthermore, if $x_{(i,\frac{j\log n}{2}+\alpha)}$
appears positive in $c$, we connect $c_2$ to all $m^i_{j'}$ such that the
binary representation of $j'$ has a $1$ in position $\alpha$. If on the other
hand $x_{(i,\frac{j\log n}{2}+\alpha)}$ appears negative in $c$, we connect
$c_2$ to all $m^i_{j'}$ such that the binary representation of $j'$ has a $0$
in position $\alpha$. In other words, we connect $c_2$ to all the middle
vertices to which $\ell^i_j$ could be connected and are consistent with an
assignment that satisfies $c$ using the current literal. After repeating the
above for all literals of each clause the construction is complete. We set the
target cost to be $B=2m+12L\log n$.

Before we argue about the correctness of the reduction, let us observe that if
the reduction preserves the satisfiability of $\phi$, then we obtain the
theorem, because the instance we constructed has vertex cover $k=O(n/\log n)$
and size polynomial in the size of $\phi$.  Indeed, as we argued the choice
gadgets have $O(n/\log n)$ vertices in total, and all further edges we added
have an endpoint in a choice gadget. If there was an algorithm solving the new
instance in time $k^{o(k)}n^{O(1)}$, this would give a $2^{o(n)}$ algorithm to
decide~$\phi$.

Regarding correctness, let us first observe that if $\phi$ is satisfiable, we
can obtain a valid solution using the intuitive translation from assignments to
choice gadget solutions we gave above. In particular, for each
$i\in\{1,\ldots,3\log n\}$ and $j\in\{0,\ldots,2L-1\}$, we consider the
assignment to variables $\{x_{(i,\frac{j\log n}{2})},\ldots,
x_{(i,\frac{(j+1)\log n}{2}-1)}\}$ as a binary number, which must have a value
$j'$ between $0$ and $\sqrt{n}-1$. We then connect both $\ell^i_j, r^i_j$ to
$m^i_{j'}$. Repeating this satisfies all demands internal to choice gadgets and
uses $3\log n \cdot 4L = 12 L \log n$ edges. Consider now a clause $c$ and the
demand from $c_1$ to $c_2$. Since we started with a satisfying assignment, $c$
must contain a true literal, say involving the variable $x_{(i,\frac{j\log
n}{2}+\alpha)}$. We select the edge from $c_1$ to $\ell^i_{j}$. Furthermore, we
observe that $c_2$ must be a neighbor of all vertices $m^i_{j'}$ such that the
bit in position $\alpha$ of the binary representation of $j'$ agrees with the
value of $x_{(i,\frac{j\log n}{2}+\alpha)}$. Since $\ell^i_{j}$ is already
connected to such a $m^i_{j'}$, we select the edge from that vertex to $c_2$ to
satisfy the demand for this clause. We have therefore spent $2m$ further edges
for the clause demands and have used a budget of exactly $B$.

For the converse direction, suppose we have a solution of cost $B$. We first
observe that each vertex $r^i_j$ must be connected to a middle vertex
$m^i_{j'}$, since all right vertices are terminals, but such vertices only have
edges connecting them to middle vertices. Recall that, for each $i,j$, the left
vertex $\ell^i_j$ must be in the same component of the solution as $r^i_j$,
since there is a demand between these two vertices. Hence, each~$\ell^i_j$ is
in the same component of the solution as some $m^i_{j'}$. We now slightly edit
the solution as follows: suppose there exists a vertex $\ell^i_j$ which is not
directly connected in the solution to any middle vertex~$m^i_{j'}$. Since this
vertex is in the same component as one such vertex $m^i_{j'}$, we add to the
solution the edge connecting them, and since this creates a cycle, remove from
the solution another edge incident on $\ell^i_j$. Doing this repeatedly ensures
that each $\ell^i_j$ is connected to a middle vertex $m^i_{j'}$ in the solution
without increasing the total cost.

We now observe that since each $\ell^i_j$ and each $r^i_j$ is connected to at
least one middle vertex $m^i_{j'}$ in the solution, this already uses a cost of
$3\log n\cdot 4L = 12 L\log n$. Furthermore, for each clause we have
constructed two terminals, each of which must use at least one of its incident
edges, giving an extra cost of~$2m$. Since our budget is exactly $2m+12L\log
n$, we conclude that each terminal constructed for a clause is incident on
exactly one edge, and each $\ell^i_j$ and each $r^i_j$ is connected to exactly
one middle vertex. Crucially, these observations imply the following fact: if
for some $i,j,j'$ we have that $\ell^i_j$ and $m^i_{j'}$ are in the same
component of the solution, then the edge connecting $\ell^i_j$ and $m^i_{j'}$
is part of the solution. To see this, observe that any path connecting
$\ell^i_j$ and $m^i_{j'}$ that is not a direct edge would need to have length at
least $3$. However, no clause terminal can be an internal vertex of such a
path, since clause terminals have degree $1$ in the solution. Furthermore, if
we remove clause terminals from the graph, left and right vertices also have
degree $1$ in the remaining solution, so such vertices also cannot be internal
in the path. Finally, middle vertices are an independent set, so it is
impossible for all internal vertices of a path of length at least $3$ to be
middle vertices.

Armed with the observation that $\ell^i_j$ and $m^i_{j'}$ are in the same
connected component of the solution if and only if they are directly connected,
we are ready to extract a satisfying assignment from the Steiner forest. For
each $i,j$, if $\ell^i_j$ is connected to $m^i_{j'}$ we write $j'$ in binary
and assign to variable $x_{(i,\frac{j\log n}{2}+\alpha)}$, for
$\alpha\in\{0,\ldots,\frac{log n}{2}-1\}$ the value in  position $\alpha$ of
the binary representation of $j'$. We claim that this assignment must be
satisfying. Indeed, consider the clause $c$, and the terminals $c_1,c_2$ which
represent it. Since these terminals have a demand, they must be in the same
component. Because $c_1$ has at most three neighbors which are in different
choice gadgets (as each clause contains variables from distinct groups), we can
see that $c_1$ must be connected to some $\ell^i_j$ and $c_2$ to some~$m^i_{j'}$
in the solution, such that $\ell^i_j$ and~$m^i_{j'}$ are in the same
component, and are therefore directly connected.  But if $\ell^i_j$ is directly
connected to~$m^i_{j'}$ this means that the assignment we extracted from
$\ell^i_j$ gives a value to a variable $x_{(i,\frac{j\log n}{2}+\alpha)}$ which
satisfies the clause $c$, hence we have a satisfying assignment.  \end{proof}


